\section{Atomic gauges and finite-coordinate polars}
\label{presentation:constructions}

The forest space also admits two direct constructions. Identify a
coefficient vector $x\in c_{00}(\mathcal F)$ with
$\sum_v x(v)e_v^*$, and let $\mathcal A$ be the atoms defined above,
including zero. Define
\[
 a(x)=\inf\{\lambda>0:x\in\lambda\operatorname{conv}\mathcal A\},
 \qquad
 q(x)=\sup_{\substack{u\in c_{00}(\mathcal F)\\\rho(u)\le1}}
       \left|\sum_v x(v)u(v)\right|.
\]
Let $X_a$ and $X_q$ be the respective completions.

\begin{proposition}\label{presentation:identity}
The coordinate identity satisfies $a(x)=q(x)=\|x\|$ on
$c_{00}(\mathcal F)$ and extends to surjective linear isometries
\[
 X_a=X_q=X=J^*.
\]
The finite unit balls and completed atomic ball are therefore exactly
those of Lemma~\ref{lem:finite-hull} and
Corollary~\ref{cor:closed-atomic-ball}.
\end{proposition}

\begin{proof}
The convex hull of $\mathcal A$ is symmetric and absorbing. Disjointness
bounds every atom's coordinates by one, while singleton segments make
$x/\|x\|_2$ an atom for $x\ne0$. Hence its gauge is a norm and
\begin{equation}\label{presentation:gauge-bounds}
 \|x\|_\infty\le a(x)\le\|x\|_2.
\end{equation}
For every $\lambda>a(x)$ there is a finite convex representation
\[
 x=\sum_{k=1}^N\theta_k A_k,\qquad
 \theta_k\ge0,\quad\sum_k\theta_k=1,
 \qquad
 A_k=\sum_i c_{ki}d_{S_{ki}},\quad
 \sum_i c_{ki}^2\le\lambda^2.
\]
Indeed, the defining infimum gives such a representation at some radius
less than $\lambda$; adjoining the zero atom permits any larger radius.
No attainment of the gauge infimum is used here.

If every segment meets $E\subset\mathcal F$ in a segment or the empty
set, then $P_E$ preserves atoms: intersect their disjoint segment
families with $E$ and discard empty terms. Thus $P_E$ is contractive for
$a$ and extends to $X_a$. This applies to ancestral sets, descendant
subtrees, and unions of incomparable descendant subtrees. Along a finite
ancestral exhaustion, $P_Hx\to x$ in $X_a$, because these projections
are contractive and eventually fix each finite coordinate vector.

Density of $c_{00}(\mathcal F)$ in $J$ gives $q(x)=\|x\|$; tests
approaching the unit ball may be rescaled to have norm at most one.
Every atom belongs to $B_X$, so $\|x\|\le a(x)$. Conversely, place the
support of a nonzero $x$ in a finite ancestral set $H$.
Lemma~\ref{lem:finite-hull} expresses $x/\|x\|$ as a finite convex
combination of atoms supported in $H$, giving $a(x)\le\|x\|$.
The zero vector is immediate. This proves equality of the three norms.
Coordinate density in Lemma~\ref{lem:forest-structure} makes both
completion maps onto $X$.

In particular the component duality already proved in that lemma
applies to the finite polar construction without a change of norm:
$q_h=q|_{c_{00}(T_h)}$ is the norm inherited from $J_h^*$, and
\[
 q(x)^2=\sum_h q_h(x|_{T_h})^2
 \qquad(x\in c_{00}(\mathcal F)).
\]
On each component the inherited bounds are
\[
 \frac{\|x\|_2}{\sqrt{h+1}}\le q_h(x)\le\|x\|_2
 \qquad(x\in c_{00}(T_h)).
\]
Thus the polar completion has the asserted component decomposition
and is reflexive.
\end{proof}

\begin{remark}[Rescaled finite approximants]\label{presentation:rescaled-closed-ball}
There is also a direct closed-ball argument from the strict gauge
infimum. If $a(x)\le1$ in the completion, choose finite $x_n$ with
$a(x_n-x)<1/n$. Then $a(x_n)\le1+1/n<1+2/n$, so the strict
infimum and balance of the convex hull give
$x_n/(1+2/n)\in\operatorname{conv}\mathcal A$. These vectors converge
to $x$. Conversely, every atom has norm at most one, so the closed
convex hull is contained in the completed unit ball. The same argument
applies to the infinite-tree gauge of Lemma~\ref{binf:finite-hull}.
\end{remark}

The atomic gauge also gives a direct proof of the subtree formula.

\begin{proposition}\label{presentation:subtrees}
Let $(r_j)$ be a finite or countable family of incomparable vertices,
and put $U_j=\{v:r_j\preceq v\}$. For $x\in X_a$ supported in
$\bigcup_jU_j$,
\[
 a(x)^2=\sum_j a(P_{U_j}x)^2.
\]
\end{proposition}

\begin{proof}
First take a finite family and $x\in c_{00}(\mathcal F)$. A segment
meets at most one $U_j$. Consequently the squared coefficient norms
of an atom's restrictions sum to at most one. The seminorm
$z\mapsto(\sum_j a(P_{U_j}z)^2)^{1/2}$ is at most one on every atom,
so it is at most $a(z)$ by convexity and passage through radii
$\lambda>a(z)$.

For the reverse inequality, choose $\lambda_j>a(P_{U_j}x)$.
Write each $P_{U_j}x$ as a finite convex average
$\sum_{k_j}\theta_{j,k_j}A_{j,k_j}$ of disjoint-segment sums with
coefficient norm at most $\lambda_j$. Project these representations
onto $U_j$, preserving that bound. For each tuple $(k_1,\ldots,k_m)$,
the sum $\sum_j A_{j,k_j}$ still uses disjoint segments and has
coefficient norm at most $(\sum_j\lambda_j^2)^{1/2}$. Give this sum
the product weight $\prod_j\theta_{j,k_j}$. These weights are
nonnegative and sum to one, and their average is
$\sum_jP_{U_j}x=x$. Therefore
\[
 a(x)\le\left(\sum_j\lambda_j^2\right)^{1/2}.
\]
Let each $\lambda_j$ decrease to $a(P_{U_j}x)$, including when this
value is zero. This proves the finite-coordinate identity.

Projection onto the union of the subtrees preserves atoms. Apply it
to finite-coordinate approximants of $x$ to obtain the finite-family
identity in the completion. For a countable family, the finite-subtree
truncations converge: a finite-coordinate approximant in the union
meets only finitely many subtrees, and projection onto the remaining
subtrees has norm at most one. The finite-family formula passes to
the limit and proves the asserted sum.

Taking the component roots gives the atomic construction's outer
$\ell_2$ sum directly. Project atomic representations onto $T_h$.
Every resulting atom has Euclidean norm at most $\sqrt{h+1}$, since
its segments are disjoint and each has at most $h+1$ vertices. Convexity and the strict-radius limit therefore give
$\|x\|_2\le\sqrt{h+1}\,a(x)$; the reverse Hilbert bound is
\eqref{presentation:gauge-bounds}. The components are thus
Hilbert-isomorphic, and the $\ell_2$-sum reflexivity argument of
Lemma~\ref{lem:forest-structure} applies to this direct construction.
\end{proof}

\begin{remark}\label{presentation:weak-null-coordinates}
The gauge gives a direct whole-dual proof of coordinate weak nullity.
For $L\in X_a^*$, \eqref{presentation:gauge-bounds} yields
\[
 |L(x)|\le\|L\|a(x)\le\|L\|\|x\|_2
 \qquad(x\in c_{00}(\mathcal F)).
\]
Thus $L$ extends continuously to $\ell_2(\mathcal F)$ and
$(L(e_v^*))_{v\in\mathcal F}$ is square summable. Its values tend to
zero along every sequence of distinct vertices. In particular, the
successive child increments are weakly null in $X_a$.
\end{remark}

\begin{remark}\label{presentation:infinite-height}
The full-dual identification uses the finite component heights. On
$T_\infty=\mathbb N^{<\omega}$, let $J_\infty$ be the completion
under the same segment norm, and let $B_\infty$ be the norm-closed
coordinate span in $J_\infty^*$, as in Section~\ref{binf:definition}.
Then
\[
 B_\infty\subsetneq J_\infty^*.
\]
Indeed, fix an infinite branch $\beta$. The functional
$f_\beta(u)=\sum_{v\in\beta}u(v)$ on finite vectors has norm one:
its value is a finite segment sum, and evaluation on any coordinate
unit vector of the branch attains one. It therefore extends to
$J_\infty$. If $y$ is a finite coordinate functional, choose a vertex
$v\in\beta$ beyond its support. Since $\rho(e_v)=1$,
\[
 \|f_\beta-y\|\ge|(f_\beta-y)(e_v)|=1.
\]
The same lower bound holds on the closed coordinate span $B_\infty$.
The zero functional gives the reverse bound, so
$\operatorname{dist}(f_\beta,B_\infty)=1$. Thus completing the finite
coordinate polar norm yields the proper coordinate subspace here.
\end{remark}
